\appendix
\raggedbottom
\makeatletter
\renewcommand{\theHfigure}{appendix.\thefigure}
\renewcommand{\theHtable}{appendix.\thetable}
\makeatother
\section{Ethical considerations}

We simulate criminal sentencing, a setting where decisions affect defendants' liberty. Every case is
procedurally generated and fictional, and every prompt says so. This study does not evaluate language
models as sentencing tools. With the guideline removed, model output responds to the arbitrary
forecast displacement ($\pullToolN$). The sentences are measurements of model behavior on synthetic
cases, not judgments about real defendants.

\section{The memory experiment}
\label{app:memory}

Deployed agent architectures more often show an agent a summary of itself than a summary of
its neighbours, so we ran a second experiment on the same cases. Agents in the
own-record arm receive a profile of their past decisions, refreshed every four
cases, giving their mean sentence, their average departure from the advisory range, how their
sentences differed by remorse and by assistance, and their most frequent rationale. Agents in
the no-memory arm decide the identical cases in the identical order with a fresh
context each time, and four baseline cases return at the end in paraphrased form, which gives
a within-agent test that no change in case composition can produce.

Without memory, the estimated trend across sixteen cases is $\driftNo$ (SE \driftNose). With an
own-record summary, the estimated trend is $\driftOwn$ (SE \driftOwnse), which is not distinguishable
from zero. The paired difference between the arms on the same case at the same position is
$\armgap$ (SE \armgapse). On returning cases, the estimated shifts are $\reexNo$ without memory
and $\reexOwn$ with memory.

We tested own-record summaries because they provide numerical information similar to peer-supplied
sentences. Each arm contains three agents and sixty decisions, limiting precision for smaller effects.

\section{What was collected}
\label{app:census}

The original arms were drawn on 17--18 August 2026 at default decoding settings, one draw per
cell, using only the prompts in the artifact repository's \texttt{experiment/prompts.txt} file; staged collection and later agent additions
explain the unequal arm sizes.

The registered cross-model arms were collected from 30 September to 1 October through the released
harness, which creates a fresh context for each prompt and inherits decoding settings. Parse failures
were recorded without redraws; one Haiku structure-matched cell failed, leaving \blHaiMSN\ rather than
\blHaiMBN\ decisions. The GPT-6 Luna sample used 26 fresh contexts across three arms and 16 cases
(1,248 decisions). It used isolated contexts with model \texttt{gpt-6-luna} and extra-high
reasoning; all replies parsed, and no cells were redrawn.

For the sentencing study, a GPT-6 Sol run was collected on 4 October with the same 1,248 cells, prompt text, cases, displayed
numbers, and displacement assignments. It used model \texttt{gpt-6-sol} at high reasoning effort.
Six first attempts returned no model output and were retried once under the protocol; all final
records parsed, and no tools were called. The Luna run used extra-high reasoning, so that setting
differs between the two model samples.

\section{Design details}
\label{app:design}

Crossing offense levels 16 and 24 with criminal history categories I and IV yields the four
advisory ranges used throughout: 21--27, 33--41, 51--63, and 77--96 months (Figure~\ref{fig:design}).

\begin{figure}[H]
\centering
\includegraphics[width=0.80\textwidth]{../figures/fig_design.pdf}
\caption{(a) Anchor values at four displacements around the advisory midpoint. (b) Identical
numbers under peer and forecast labels; their slope difference is the peer premium.}
\label{fig:design}
\end{figure}

Write $F_c \in \{0,1\}^4$ for case $c$'s factor vector and partition the sixteen design points
into four groups $g_1,\dots,g_4$ subject to
\begin{equation}
\sum_{c \in g_k} F_{cf} \;=\; \tfrac{1}{2}\,|g_k| \qquad \text{for every group } k
\text{ and every factor } f,
\label{eq:balance}
\end{equation}
which a backtracking search satisfies exactly at $|g_k| = 4$. Agent $i$ then receives
\begin{equation}
\delta_{ic} \;=\; d_{\,(k(c) + i) \bmod 4}, \qquad
d \;=\; (-0.30,\,-0.15,\,+0.15,\,+0.30),
\label{eq:alloc}
\end{equation}
where $k(c)$ indexes the group containing $c$. Every agent sees each displacement exactly four
times, and \eqref{eq:balance} makes the four factors split evenly within each displacement level.

\section{Inference}
\label{app:inference}

Drawing $\zeta_{ic} \sim \mathrm{Unif}\{0,1\}$ independently per cell and relabelling
$P^{(b)}_{ic} = \zeta_{ic} P_{ic} + (1-\zeta_{ic})(1 - P_{ic})$ generates the exact null
distribution of $\hat\Delta$, from which
\begin{equation}
p \;=\; \frac{1 + \#\{b : |\hat\Delta^{(b)}| \ge |\hat\Delta|\}}{1 + B},
\qquad
q_{.95} \;=\; \text{95th percentile of } |\hat\Delta^{(b)}|,
\label{eq:ri}
\end{equation}
with $B = 9{,}999$.

We quote all three clusterings for the premium contrasts and the case-clustered figure for per-arm
pull. Cluster-robust $p$-values use a wild cluster bootstrap with Webb six-point weights
\citep{webb2023reworking,cameron2008bootstrap,roodman2019fast}, which we use because
agent-clustered inference rests on six clusters. Per-arm pull uses a second randomization null,
permuting the displacement within each agent across the design's own factor-balanced case groups;
\eqref{eq:ri} describes the attribution-label null used for premiums. Population agreement is
summarised by the between-agent spread on a common case,
\begin{equation}
\bar{\sigma} \;=\; \frac{1}{|\mathcal{C}|}\sum_{c \in \mathcal{C}}
\operatorname{sd}_{i}\!\left( s_{ic} / m_c \right),
\label{eq:spread}
\end{equation}
the mean over cases of the standard deviation across agents of the sentence expressed in units of
the case midpoint.

Across \yokedCells\ cells in which the same agent saw the same case under more than one
attribution, the displayed numbers were identical in every paired cell. The four
case factors remain balanced across displacement levels, with the largest departure from an even
split equal to \worstImbalance. Adding offense-type fixed effects leaves the guideline-present
premium at $\rbGType$ (SE \rbGTypese).

Adding a guideline indicator $G$ with $\delta \times G$, $P \times G$ and $\delta \times P \times
G$ terms gives a moderation model whose triple interaction is $\triple$ (SE \triplese,
$p=\triplep$). Only arms with the original closing lines were run under both guideline conditions,
and their peer-premium estimates are near zero; we therefore report the conditions separately.

\section{Pull by arm}
\label{app:cells}

Arm-level pulls include all arm agents; paired premiums use shared cells (Table~\ref{tab:cells}).

\begin{table}[H]
\centering
\caption{Opus 5 pull by block and source. Case-clustered SEs; randomization $p$-values. No-number rows show mean deviation.}
\label{tab:cells}
\footnotesize
\renewcommand{\arraystretch}{0.86}
\begin{tabular}{llrrrr}
\toprule
Numbers attributed to & Block variant & $n$ & Pull $\hat{\pi}$ & SE & $p$ \\
\midrule
\multicolumn{6}{l}{\emph{Guideline removed}} \\
Other judges of this bench & none & \nPeerBare & $\pullPeerBare$ & \sePeerBare & --- \\
Other judges of this bench & structure-matched & \nPMatch & $\pullPMatch$ & \sePMatch & --- \\
Other judges of this bench & paraphrase one & \nPFree & $\pullPFree$ & \sePFree & --- \\
Other judges of this bench & paraphrase two & \nPOwn & $\pullPOwn$ & \sePOwn & --- \\
Statistical forecast & none & \nToolBare & $\pullToolBare$ & \seToolBare & --- \\
Other judges of this bench & ``decide independently'' & \nPeerN & $\pullPeerN$ & \sePeerN & \riPeerNp \\
Statistical forecast & ``advisory only'' & \nToolN & $\pullToolN$ & \seToolN & \riToolNp \\
Docketing artifact & ``no legal significance'' & \nClerN & $\pullClerN$ & \seClerN & \riClerNp \\
\addlinespace[1pt]
\multicolumn{6}{l}{\emph{Guideline present}} \\
Other judges of this bench & ``decide independently'' & \nPeerG & $\pullPeerG$ & \sePeerG & \riPeerGp \\
Statistical forecast & ``advisory only'' & \nToolG & $\pullToolG$ & \seToolG & \riToolGp \\
\midrule
\multicolumn{2}{l}{No numbers shown, guideline present} & \NaGn &
  \multicolumn{3}{l}{mean deviation $\NaGdev$} \\
\multicolumn{2}{l}{No numbers shown, guideline removed} & \NaNn &
  \multicolumn{3}{l}{mean deviation $\NaNdev$} \\
\bottomrule
\end{tabular}
\end{table}

\section{Cross-model estimates}
\label{app:models}

With a guideline and no numbers, between-agent spread is \NaGsd; without the guideline it is
\dispFull. Leave-one-out estimates for the no-guideline condition range from \dispLooLo\ to
\dispLooHi, so the estimated magnitude varies with the four-agent sample.
Guideline-present premiums per model are $\xgOpus$ (SE \xgOpusse\ on case, $n=\xgOpusn$) on Opus~5,
$\xgSonnet$ (SE \xgSonnetse) on Sonnet~5 and $\xgHaiku$ (SE \xgHaikuse) on Haiku~4.5, the latter two
over \xgSonnetn\ decisions from one agent each.

We ran both displaced-anchor arms on three models of the same family, with and without the
advisory guideline. With the guideline present,
pull increases as model size decreases, from near zero on Opus 5 to roughly two-thirds on Haiku 4.5.
Without the guideline, the spread in pull across the three models is smaller, although Opus 5 remains
below full adoption.

The estimated peer premiums are $\bmOpus$ on Opus 5, $\bmHaiku$ on Haiku 4.5, and $\bmSonnet$ on
Sonnet 5. The smaller-model estimates use $\bmSonnetn$ decisions each and are too imprecise to
separate. We make no cross-model scaling claim from these three estimates.

The GPT-6 Luna and Sol premiums use the same cases and prompts. Their full estimates are reported in
Table~\ref{tab:gpt-cross}; the comparisons across models are descriptive because the reasoning
settings differ.

\begin{table}[H]
\centering
\caption{GPT-6 peer premiums relative to the bare forecast. Each contrast uses 416 observations
per arm and case-clustered 95\% intervals. The two runs used identical prompts and allocations.}
\label{tab:gpt-cross}
\scriptsize
\setlength{\tabcolsep}{3.5pt}
\renewcommand{\arraystretch}{0.95}
\begin{tabular}{lllrcrr}
\toprule
Model & Effort & Contrast & Premium [95\% CI] & $p$ & MDE & $n$ \\
\midrule
GPT-6 Luna & xhigh & Structure-matched & $\xlMSEst$ [$\xlMSLo$, $\xlMSHi$] & $\xlMSP$ & $\xlMSMde$ & \xlMSN \\
GPT-6 Luna & xhigh & Bare & $\xlMBEst$ [$\xlMBLo$, $\xlMBHi$] & $\xlMBP$ & $\xlMBMde$ & \xlMBN \\
GPT-6 Sol & high & Structure-matched & $\solMSEst$ [$\solMSLo$, $\solMSHi$] & $\solMSP$ & $\solMSMde$ & \solMSN \\
GPT-6 Sol & high & Bare & $\solMBEst$ [$\solMBLo$, $\solMBHi$] & $\solMBP$ & $\solMBMde$ & \solMBN \\
\bottomrule
\end{tabular}
\end{table}

\section{Confidence and exact matching}
\label{app:selfreport}

Confidence uses \rbConfDistinct\ of its ten values, and
\rbConfTop\% of ratings are six or seven. We treat confidence as an uncalibrated self-report and
describe it across arm means. Against the same bare forecast, the first independence paraphrase yields
$\pfPrem$ ($p=\pfPremp$), while the second yields $\poPrem$ ($p=\poPremp$)
(Section~\ref{sec:confidence}; Figure~\ref{fig:confidence}).

\begin{figure}[H]
\centering
\includegraphics[width=0.80\textwidth]{../figures/fig_confidence.pdf}
\caption{(a) Mean confidence against pull for each model-arm cell; the dashed line and correlation
($r=\confCorr$) use only the six Opus 5 arm means (arm $n=16$--416). (b) Equal-weighted mean, by
provider, of model-level shares matching one of the three displayed numbers, for the three arms common
to all five models.}
\label{fig:confidence}
\end{figure}

\section{The cascade floors}
\label{app:cascade}

Shuffling independent agents into pseudo-speaking orders yields a \plCopy\% exact-match rate and a
slope of $\plBeta$ ($[\plBetalo, \plBetahi]$). Round-number clustering and case-level variation can
produce these matches without peer influence. With case fixed effects, the same placebo returns
$\plBetaFE$ ($[\plBetaFElo, \plBetaFEhi]$), consistent with the negative sign
\citet{angrist2014perils} predicts when case dummies place predecessor outcomes on both sides of a
within-case comparison. In the cascade, the fixed-effect slope is $\cascBetaFE$ (SE \cascBetaFEse),
compared with the plain slope $\cascBeta$ (SE \cascBetase). The fixed-effect estimate is farther
from its placebo floor; the plain estimate is close to $\plBeta$. The placebo uses at most three
predecessors, compared with five in the cascade, so it does not provide a floor matched to the later
seats. The reflection problem \citep{manski1993identification} does not arise because each agent's
predecessors are fixed before it decides.

The first four agents speak in an order randomised per case; the last two were added afterwards in
fixed positions, so position effects are identified over the first four. Exact reproduction rates
at the second, third and fourth positions are \cascCopyOne\%, \cascCopyPosTwo\% and
\cascCopyPosThree\%, respectively, with sixteen decisions at each position. In the docketing arm,
which shows three numbers in the same layout, the exact-match rate is \exCD\%. The final-pair spread
is $\cascLate$, compared with $\cascGuideSd$ for the independent population with a guideline and no
peer information. The estimates use different agent sets, so we do not calculate a fold-change.

\section{Robustness of the premium}
\label{app:robust}

\subsection{Design checks}
Across \yokedCells\ cells in which the same agent saw the same case under more than one
attribution, the numbers shown differed in \yokedMismatch\ cells, so the yoke is exact. The four
case factors remain balanced across displacement levels, with the largest departure from an even
split equal to \worstImbalance. Adding offense-type fixed effects leaves the guideline-present
premium at $\rbGType$ (SE \rbGTypese).

\subsection{Sentence values and displayed numbers}
Over \rbSentN\ decisions agents use \rbDistinct\ distinct sentence values, \rbDivSix\% divisible by
six months, \rbDivTwelve\% by twelve and \rbDivFive\% by five, with \rbMode\ months the single most
common answer. Of the numbers shown to agents, \rbAncSix\% are divisible by six and \rbAncFive\% by
five, and no record shows three round numbers at once. Decisions reproduce a displayed number in
\rbExact\% of cases (\rbExactN\ decisions), of which \rbExactRound\% are divisible by six; matching
runs \rbMatchRound\% where some displayed number is round against \rbMatchPlain\% where none is.
The original-sentence comparison is not statistically distinguishable from zero with round numbers
excluded, $\rbNROS$ ($p=\rbNROSp$, $n=\rbNROSn$). Exact matching is more common when at least one
displayed number is round than when none is. We therefore compare copying with a shuffled floor.
Because both attribution arms show the same numbers, round-number matching does not change the
estimated premium.

\subsection{Premium by displacement magnitude}
We estimate the premium at each displacement. The bare premium is flat, $\rbMagMBLo$ at
$|\delta|=0.15$ and $\rbMagMBHi$ at $|\delta|=0.30$; the structure-matched premium rises from
$\rbMagMSLo$ ($p=\rbMagMSLop$) to $\rbMagMSHi$ ($p=\rbMagMSHip$). It is larger at the far than at
the near displacement. A single slope averages these displacement-specific estimates.
The original-sentence comparison is null at both magnitudes: $\rbMagOSLo$
($p=\rbMagOSLop$) and $\rbMagOSHi$ ($p=\rbMagOSHip$).

\subsection{Alternative outcome scales}
On the log-ratio scale, the matched and bare premiums are $\rbLogMS$ ($p=\rbLogMSp$) and
$\rbLogMB$ ($p\rbLogMBp$). On the within-case midrank scale, they are $\rbRankMS$
($p=\rbRankMSp$) and $\rbRankMB$ ($p\rbRankMBp$). The original-closing comparison remains
null on both scales ($\rbLogOS$, $p=\rbLogOSp$; $\rbRankOS$, $p=\rbRankOSp$). Since ranks and
ratios use different scales, we compare signs and $p$-values, not magnitudes.

\section{Behaviour under the equivalence sentence}
\label{app:equiv}

The dispersion and copying shifts in Section~\ref{sec:equiv} are post-hoc diagnostics, not
pre-registered. On Opus~5, outcome standard deviations are $0.235$ and $0.242$ in the
structure-matched and bare pairs, versus $0.336$ and $0.329$ with the equivalence and reliability
sentences. Exact displayed-number matches are $29\%$, $25\%$, and $28\%$ in the earlier arms,
versus $8\%$, $12\%$, $6\%$, and $8\%$ in the four later arms. Sentences outside the span of the
three numbers rise from $31$--$57\%$ to $80$--$84\%$.

In both peer and forecast arms, agents independently said the numbers added no evidence beyond the
case file. One wrote that they ``derive from the identical case file and add no new information
about the defendant, so I gave them no weight as evidence.'' Responses were recorded verbatim and
not re-drawn.

\section{Case factors}
\label{app:factors}

Without numbers, severity and prior record increase sentence length by \NaNsev\% and \NaNprior\%
(Figure~\ref{fig:factors}). With a guideline, \rangeInside\% of decisions fall within its range,
\rangeBelow\% below, and \rangeAbove\% above; departures are almost entirely downward.
Proportionality is the most common stated rationale with a guideline (\RatGtopv\%); deterrence is
the most common without it (\RatNtopv\%). The factor analysis estimates larger effects for severity and prior
record (\facSev\%, \facPrior\%) than for assistance and remorse ($\facCoop$\%, \facRem\%).

\begin{figure}[H]
\centering
\includegraphics[width=0.88\textwidth]{../figures/fig_factors.pdf}
\caption{Case-factor effects on sentence length by guideline condition.}
\label{fig:factors}
\end{figure}

\section{What each arm can exclude}
\label{sec:power}

Of \pwTotal\ estimates, \pwUnder\ miss the criterion; all misses are in smaller-model cells
(Table~\ref{tab:power}).

\begin{table}[H]
\centering
\caption{Minimum detectable effect (MDE) is $2.80\times\mathrm{SE}$ for 80\% power at two-sided
5\%. ``Powered'' indicates whether each estimate meets its Claude benchmark.}
\label{tab:power}
\footnotesize
\renewcommand{\arraystretch}{0.9}
\begin{tabular}{lrrrrc}
\toprule
Estimate & $n$ & Est. & SE & MDE & Powered \\
\midrule
Structure-matched premium, Opus 5 & 192 & $+0.206$ & 0.079 & 0.222 & yes \\
Bare-block premium, Opus 5 & 128 & $+0.323$ & 0.047 & 0.132 & yes \\
Cost of the peer closing sentence & 128 & $-0.279$ & 0.074 & 0.208 & yes \\
Cost of paraphrase one & 96 & $-0.326$ & 0.066 & 0.184 & yes \\
Cost of paraphrase two & 96 & $-0.198$ & 0.070 & 0.196 & yes \\
Premium with both original closings & 128 & $-0.006$ & 0.068 & 0.191 & yes \\
Premium, Opus 5, guideline present & 192 & $-0.018$ & 0.041 & 0.116 & yes \\
Premium, Opus 5, guideline removed & 128 & $-0.006$ & 0.068 & 0.191 & yes \\
Premium, Sonnet 5, guideline present & 32 & $+0.204$ & 0.056 & 0.156 & yes \\
Premium, Sonnet 5, guideline removed & 32 & $+0.065$ & 0.290 & 0.812 & \textbf{no} \\
Premium, Haiku 4.5, guideline present & 32 & $+0.211$ & 0.112 & 0.315 & yes \\
Premium, Haiku 4.5, guideline removed & 32 & $+0.071$ & 0.119 & 0.333 & \textbf{no} \\
Premium, smaller models pooled & 128 & $+0.138$ & 0.114 & 0.318 & yes \\
GPT-6 Luna, structure-matched & 832 & $+0.220$ & 0.049 & 0.136 & yes \\
GPT-6 Luna, bare blocks & 832 & $+0.256$ & 0.049 & 0.138 & yes \\
GPT-6 Sol, structure-matched & 832 & $+0.239$ & 0.054 & 0.150 & yes \\
GPT-6 Sol, bare blocks & 832 & $+0.113$ & 0.044 & 0.123 & yes \\
Extra pull from removing guideline, Opus 5 & 320 & $+0.541$ & 0.071 & 0.199 & yes \\
Extra pull from removing guideline, Sonnet 5 & 64 & $+0.632$ & 0.152 & 0.426 & \textbf{no} \\
Extra pull from removing guideline, Haiku 4.5 & 64 & $+0.357$ & 0.063 & 0.176 & yes \\
\bottomrule
\end{tabular}
\end{table}


\section{Civil-damages results across model lineages}
\label{app:damages}

We applied the displaced-number design to sixteen fictional civil claims, crossing severity,
liability, mitigation, and documentation. D1 used Claude Opus~5. Its proximity measures were
selected after examining the first eight agents, so those panel estimates are exploratory; the
measures and direction were then registered before the two-agent confirmation sample (A12 and A01)
was examined. On those 32 matched pairs, in-band awards were 14/32 under the peer label and 26/32
under the tool label ($p=0.0042$, two-sided exact sign test), awards within two percent of the
central figure 10/32 and 24/32 ($p=0.00052$), and exact central matches 2/32 and 23/32
($p=9.54\times10^{-7}$). The peer-minus-tool difference in proportional distance was $+0.140$
(95\% $t$ $[0.032,0.248]$, MDE $0.142$). The registered rule calls this suggestive; see
\texttt{damages/AMENDMENT-D1-proximity.md}.

D2 was registered before collection and used the same ten agents, claims, primary arms, and
rendered prompts as the complete D1 panel; every prompt hash matched its D1 counterpart. GPT-6 Sol
at high reasoning returned 319 of 320 planned responses parseably (the A07 step-3 tool cell gave
none and was not retried), leaving 159 matched pairs. Table~\ref{tab:damages-cross} compares all
ten D1 agents with D2; its D1 column pools the exploratory panel with the confirmation sample and
is descriptive. Differences are peer-minus-tool with 95\% claim-clustered $t$ intervals; binary
$p$-values are exact sign tests on discordant pairs, reported peer-only/tool-only.

\begin{table}[H]
\centering
\caption{Civil-damages proximity outcomes by model run.}
\label{tab:damages-cross}
\scriptsize
\setlength{\tabcolsep}{3pt}
\renewcommand{\arraystretch}{1.0}
\begin{tabular}{>{\raggedright\arraybackslash}p{0.15\textwidth}>{\raggedright\arraybackslash}p{0.38\textwidth}>{\raggedright\arraybackslash}p{0.38\textwidth}}
\toprule
Outcome & Claude Opus 5, D1 (160 pairs) & GPT-6 Sol high, D2 (159 pairs) \\
\midrule
Inside band &
Peer 21/160; tool 88/160; $\Delta=-0.419$ [$-0.557,-0.281$]; $p=2.17\times10^{-18}$; discordant peer-only/tool-only $2/69$ &
Peer 90/159; tool 88/159; $\Delta=+0.013$ [$-0.127,+0.152$]; $p=0.871$; discordant $20/18$ \\
Distance from center &
Mean peer 0.392; tool 0.156; $\Delta=+0.236$ [$+0.145,+0.328$]; MDE 0.120 &
Mean peer 0.274; tool 0.182; $\Delta=+0.093$ [$+0.011,+0.174$]; MDE 0.107 \\
Within 2\% of center &
Peer 11/160; tool 68/160; $\Delta=-0.356$ [$-0.485,-0.227$]; $p=2.08\times10^{-16}$; discordant $1/58$ &
Peer 72/159; tool 73/159; $\Delta=-0.006$ [$-0.101,+0.089$]; $p=1.000$; discordant $17/18$ \\
Equals center &
Peer 2/160; tool 49/160; $\Delta=-0.294$ [$-0.390,-0.198$]; $p=1.42\times10^{-14}$; discordant $0/47$ &
Peer 34/159; tool 43/159; $\Delta=-0.057$ [$-0.152,+0.039$]; $p=0.200$; discordant $15/24$ \\
\bottomrule
\end{tabular}
\end{table}

The registered D2 primary outcome does not reproduce D1's large in-band difference: the contrast is
small, reverses direction, and is not statistically decisive. This is not evidence that the arms
are equivalent. For central distance, the D2 interval excludes zero, but its MDE exceeds the point
estimate, so the registered rule calls the result suggestive rather than decisive. Neither binary
secondary outcome rejects. The strong D1 proximity pattern therefore does not generalize unchanged
to this GPT-6 Sol sample.

For continuity, the pooled displacement slope is $+0.383$ on D1 (95\% $z$ interval
$[0.185,0.581]$, MDE $0.283$, within-agent permutation $p=0.0014$) and $+0.559$ on D2
(95\% $z$ interval $[0.344,0.774]$, MDE $0.307$, $p=0.0002$). The tool-minus-peer slope premium
is $+0.215$ on D1 (MDE $0.311$, $p=0.0712$) and $+0.232$ on D2 (MDE $0.292$, $p=0.0341$); both
estimates are below their MDEs. These secondary summaries do not replace the registered proximity
outcomes. D1 used Claude subagents, whereas D2 used Codex CLI, so the cross-run difference is
descriptive across both model and execution route. The results do not separate trust in a source
label from adoption of a supplied number.
